%%%%%%%%%%%%%%%%%%%%%%%%%%%%%%%%%%%%%%%%%%%%%%%%%%%%%%%%%%%%%%%%%%%%%%%%%%%%%%%%
% BLOQUE 5 — DESARROLLO Y ANÁLISIS TÉCNICO
% Estructura obligatoria del informe TI-05 (ver ficha, cruzada con las
% pautas generales del curso). Cubre los 6 subtemas obligatorios de la
% ficha TI-05, todos en ESTE ÚNICO ARCHIVO (por preferencia del equipo,
% sin subcarpetas ni \input separados por subtema).
%%%%%%%%%%%%%%%%%%%%%%%%%%%%%%%%%%%%%%%%%%%%%%%%%%%%%%%%%%%%%%%%%%%%%%%%%%%%%%%%

\section{Desarrollo y análisis técnico}

% --- Subtema 1/6 (ficha TI-05) ---
\subsection{Modelos de ejecución y aislamiento}

% TODO (equipo): cubrir funciones como servicio (FaaS), contenedores sin
% servidor y entornos de borde; y los mecanismos de aislamiento V8,
% microVM y contenedor.

\subsubsection{Alternativas adicionales identificadas por el grupo}

% OBLIGATORIO (punto 4 de las indicaciones): la ficha lista, por
% categoría, algunos proveedores/productos y cierra con "y otros que el
% grupo debe identificar". Para cada categoría comparada aquí (FaaS,
% contenedores sin servidor, entornos de borde), agregar al menos DOS
% alternativas adicionales no mencionadas en la ficha, con su
% justificación; o, si tras la búsqueda no se encuentran más,
% declararlo expresamente y documentar la búsqueda (fuentes
% consultadas, criterios de descarte, fecha).

% --- Subtema 2/6 (ficha TI-05) ---
\subsection{Arranque en frío}

% TODO (equipo): causas del arranque en frío, cómo se mide, y
% mitigaciones (concurrencia aprovisionada, instantáneas/snapshots,
% entornos livianos). Ver también el Entregable específico de PoC/
% mediciones de arranque en frío (bloque 6).

% --- Subtema 3/6 (ficha TI-05) ---
\subsection{Facturación por consumo y punto de equilibrio}

% TODO (equipo): unidades de facturación (GB-segundo, vCPU-segundo,
% tiempo de CPU frente a tiempo de reloj, número de invocaciones,
% salida de datos) y el concepto de punto de equilibrio frente a
% instancias reservadas. El cálculo numérico del punto de equilibrio va
% en el Entregable específico "Modelo de costos" (bloque 6); esta
% subsección explica los conceptos y variables que ese modelo usa.

% --- Subtema 4/6 (ficha TI-05) ---
\subsection{Diseño orientado a eventos}

% TODO (equipo): colas, publicación/suscripción, idempotencia, entrega
% al menos una vez ("at-least-once"), y orquestación frente a
% coreografía.

% --- Subtema 5/6 (ficha TI-05) ---
\subsection{WebAssembly y WASI como entorno de borde}

% TODO (equipo): portabilidad, aislamiento y arranque submilisegundo de
% WebAssembly/WASI como entorno de ejecución en el borde.

% --- Subtema 6/6 (ficha TI-05) ---
\subsection{Límites técnicos y dependencia del proveedor (vendor lock-in)}

% TODO (equipo): tiempos máximos de ejecución, tamaño de carga útil,
% manejo de estado y estrategias de portabilidad entre proveedores. Ver
% también el cuadro comparativo de plataformas (bloque 6).

\clearpage
